\section{Lesson 1: Geography Is Destiny for Single-Member Districts}
\label{sec:lesson1}

The single most important finding from the E-track research program is also
its most sobering: for any system that uses geographically compact single-member
districts, Democratic underrepresentation relative to national vote share is a
structural inevitability, not a drawing artifact.  This lesson is established
most directly by Papers E.4 and E.5, which approach the sorting problem from
opposite directions.

\subsection{The Rodden Mechanism}

\citet{rodden2019why} identifies the mechanism precisely.  Democratic voters
are geographically concentrated in dense urban cores in ways that produce
``wasted votes'' under winner-take-all single-member elections.  A district
centered on central Detroit, Philadelphia, or San Francisco will vote 85--15
Democratic regardless of how its boundaries are drawn; the 85\% who vote
Democratic ``waste'' their votes in the sense that the district would have
produced a Democratic representative with 51\% of the vote, and the additional
34 percentage points of Democratic margin are electorally inert.  Republican
voters, by contrast, are spread across suburban and rural geographies where
their margins are more modest --- close enough to win, but not so large as
to generate systematic waste.

The result is that in a system with perfect population equality, Democratic
voters systematically produce fewer Democratic seats per Democratic vote than
Republican voters produce Republican seats.  \citet{chen2013unintentional}
showed that this effect is large enough to produce Republican seat majorities
in many states even when district lines are drawn by nonpartisan algorithms
with no partisan intent.  Our 50-state algorithmic sweep confirms this finding
with more recent data and higher geographic resolution \citep{dellaLibera2026mec}.

\subsection{Evidence from E.4: Partisan-Similarity Clustering}

Paper E.4 provides the sharpest demonstration that sorting is a structural
rather than drawing artifact \citep{dellaLibera2026e4}.  By varying the
partisan-similarity edge weight from zero (pure compactness, the baseline) to
a maximum that drives all partisan-similar tracts into the same district,
E.4 shows that the aggregate seat--vote relationship remains almost unchanged.

Specifically, when we weight tract adjacencies by partisan similarity and
optimize for partisan homogeneity within districts, the resulting maps have
95\% safe seats (versus 52\% in the baseline) and a much wider distribution
of partisan margins.  But the mean seat shares for both parties are nearly
identical to the baseline: the algorithm has changed which seats each party
wins safely and which it wins narrowly, but it has not changed how many seats
each party wins in expectation.  The seat bonus for Republicans --- the gap
between Republican seat share and Republican vote share --- is $+4.2$
percentage points in the baseline and $+4.0$ percentage points in the
partisan-similarity extreme, a difference of 0.2~pp that is not statistically
distinguishable from zero ($p = 0.41$).

This null result for partisan redistribution within a single-member geographic
system is the most direct evidence that sorting, not drawing, is the
operative constraint.

\subsection{Evidence from E.5: Partisan-Fairness Allocation}

Paper E.5 approaches the same problem from the other direction: instead of
asking what happens when we optimize for partisan homogeneity, it asks what
the algorithm must do to achieve proportionality \citep{dellaLibera2026e5}.
The answer is that it must draw a specific class of non-compact districts ---
those that crack urban Democratic concentrations into multiple surrounding
suburban districts, reducing the margin wasted in each.

The quantity and severity of these non-compact districts is a direct measure
of the geometric cost of escaping sorting.  In states with severe Democratic
urban concentration (Illinois, New York, Pennsylvania, Michigan), achieving
proportionality requires districts that would fail any reasonable compactness
threshold.  In states with more evenly distributed partisanship (Iowa,
Wisconsin, Arizona), the compactness cost of proportionality is modest.

This cross-state variation in the geographic cost of proportionality is itself
informative.  It implies that state-by-state adoption of proportionality
standards would have uneven effects: states with less sorting would achieve
proportionality at low compactness cost, while states with high urban
concentration would face a much steeper trade-off.  A national proportionality
standard would thus impose differential burdens on different states.

\subsection{The Implication: Drawing Methods Are Not the Remedy}

The combined evidence from E.4 and E.5 establishes that improving the drawing
method --- whether by algorithmic means, independent commission, or court
supervision --- cannot eliminate the seat--vote gap arising from geographic
sorting.  The gap exists at the map level, not the drawing-process level.
Courts that focus on drawing \emph{process} (whether legislators had partisan
intent) rather than drawing \emph{outcomes} (whether the resulting map is
proportional) are, by this analysis, addressing the wrong question for the
wrong reason.

This finding does not imply that drawing-process reforms are without value.
Algorithmic redistricting and independent commissions do eliminate the
\emph{additional} seat bias that partisan gerrymanders layer on top of the
structural sorting bias.  In some states, that additional bias is large
(North Carolina, Ohio, Wisconsin in the 2010s) and its elimination is
consequential.  But in states where the structural sorting bias already
dominates, no process reform changes outcomes materially.
